\section{地域、编户与帝国信息：西汉社会的多重尺度}
\label{sec:western-han-regions-households-information}

西汉的统一并没有立刻造出一片由同一套赋役、语言和市场无差别覆盖的腹地。刘邦先入关中、秦王子婴降于轵道，能够说明关中政权的军事终结，却不能说明秦故吏、县廷文书、仓储和乡里关系在何时、以何种方式被汉廷接收。关中是皇帝能够较直接调度的核心，诸侯国、边郡、旧楚地和南方新附地区则各有不同的地方精英、道路和物资来源。将“汉承秦制”理解为一项已经完成的社会事实，反而遮住了新政权必须逐县重建征发、司法和信息渠道的工作。

\subsection{轻徭薄赋是登记问题，而非一个社会状态}

前一七八年文帝减田租为三十税一，诏令把战后恢复、编户稳定和皇帝德政联结起来。它确实为观察中央如何表述负担提供清楚的锚点；但田租不是农家面对的唯一成本。劳役、转输、地方性征解、欠赋和土地占有的差异，并不会因为诏文使用同一税率而消失。没有成套的县乡簿籍，我们不能由《汉书》中的制度语言计算某一地区或某一类家庭的实收负担。应当保留的结论是：朝廷把可登记的田地和人户视作恢复国家能力的条件，而地方社会是否、如何进入这种登记，仍有待不同材料逐处检验。

前一七七年前后弛山泽禁并许可民间铸钱，又揭示另一种并不等同于田租的治理路径。矿产、钱范、技术、诸侯国财力与商人网络，使货币流通穿过郡县边界。民间和王国铸钱扩大，币质与地方财力问题随之进入中央视野；这并不意味着市场从此不受国家影响，也不说明每个农户都以钱币结算日常生活。钱范、钱币窖藏和传世食货叙事可以互证铸币和整顿的存在，却不能把窖藏数量直接换算为某州的交易总量。货币史的价值，正在于使人看见国家征收之外仍有多种资源与中介在组织社会。

\subsection{王国改革与郡县的不同接口}

七国之乱后，王国官属和军政空间被进一步收缩。周亚夫在昌邑一线断绝吴楚粮道的胜利，显示叛军同样依赖道路、粮食和沿线地方配合；它不应被化为“中央战胜地方”的抽象公式。王国没有立即消失，宗室封号、食邑和地方关系仍是汉廷必须处理的政治资源。前一四四年前后对王国制度的调整，是让任官、司法和军政更能被中枢核验的尝试，而不是消除地方社会的全部自主性。

前一三四年令郡国举孝廉，也应放进同一脉络理解。制度将地方评价接到官员任用上，使郡国能够把被认为适合仕宦的人送入中央；它既不是后来完整科举的前身，也不能证明地方声望被国家单向吸收。推荐所依赖的家族、师友、郡守和地方舆论本身就是社会关系。朝廷获得的是一条可书写、可考课的入口，地方则保留了影响谁能被看见、谁能被证明“孝”或“廉”的空间。

\subsection{河西并非一条直线向西延伸}

前一二七年卫青收复河南地，前一二一年浑邪王率部降汉，前一二〇年至前一一一年间武威、酒泉、张掖、敦煌陆续设置，构成河西经营的年序。每个节点都包含安置、屯戍、邮传、畜牧、粮食和使节联系；把它们画成汉军越过某条界线即可“占有”的过程，会掩盖道路和绿洲之间持续的供给问题。降附部众与新设郡县不是两种可以完全分开的实体，中央分类、地方自称和实际居处都可能变化。居延、肩水金关、悬泉置等文书能够让吏员、军粮、出入关和邮传的日常从帝纪背后浮现，却只是特定地点与机构保存下来的档案。

前一一一年入番禺、前一〇九年入滇、前一〇八年灭卫满朝鲜，也把港口、山地、地方首领和东北道路纳入汉廷的战略。胜负、郡县名目和大致年序可以由《史记》《汉书》对读；当地精英如何重新安排土地、劳力和贸易，未必被中央叙事连续保存。所谓“设郡”是任命与文书上的动作，不是地方社会在同一天变成关中式县乡。因而，西汉的扩张应当理解为不同地区被接入军事、财政和名义网络的速度不同，而不是一张一次完成的行政地图。

\subsection{财政集中如何进入乡里}

前一一九年五铢、盐铁专营与前一一八年前后的均输、平准，把钱币、盐铁和货物流通更紧地接入中央调度。它们回应远征、边郡与都城供应的压力，也让地方运输和官商关系变得更紧张。前九八年前后继续整顿币制，显示中央并未取得一次措施便永远稳定的货币控制；前八九年轮台诏拒绝继续在西域大规模屯田远征，则承认徭役、军费和远距离供给的承受力有限。前八一年盐铁会议保留的，是围绕专营、民生与边防成本的可说分歧，并非一张能证明全国意见的社会调查。

前六一年至前六〇年赵充国在羌地倡屯田、招抚、分化，进一步揭露财政、地方社会和边防不能分开。屯田需要土地、劳力和水源，也需要能与不同部众、边郡官员和军队沟通的安排。正史所称“羌”首先是战争和行政中的分类，不应替代各部的政治目标与社会组成。赵充国传能够保存政策争论，却不能为每一处屯田计算收成、安置人数或服从程度。西汉国家能力的边界正存在于这里：朝廷可以不断写下税率、钱制和安边方案，但这些方案总要经由地方土地、人群和道路才可能成为现实。

\subsection{来源与材料边界}

本节的年序主干取自《史记》高祖、本纪与〈平准书〉、《汉书》帝纪、〈食货志〉、〈地理志〉、匈奴、西域和赵充国相关传记，并以《盐铁论》、河西简牍、钱范、钱币窖藏和城址材料校看不同层面。帝纪适合确定诏令和大事，食货志与论辩保存制度语言，简牍和遗存照见局部执行；它们都不能单独证明全国人口、市场规模或普通家庭的动机。这样区分材料，不是削弱西汉的统一，而是使统一被理解为不断处理地区差异的一组实践。
